\documentclass[11pt,a4paper]{article}

% ---------- PAQUETES ----------
\usepackage[utf8]{inputenc}
\usepackage[T1]{fontenc}
\usepackage[spanish,es-noshorthands]{babel}
\usepackage{lmodern}
\usepackage{geometry}
\usepackage{microtype}
\usepackage{booktabs}
\usepackage{tabularx}
\usepackage{float}
\usepackage{xcolor}
\usepackage{hyperref}
\usepackage{fancyhdr}
\usepackage{enumitem}
\usepackage{listings}

% ---------- ESTILO ----------
\geometry{margin=2.4cm}

\definecolor{azul}{RGB}{19,53,92}
\definecolor{gris}{RGB}{85,95,110}

\hypersetup{
    colorlinks=true,
    linkcolor=azul,
    urlcolor=azul,
    pdftitle={Actividad 12 - Informe tecnico},
    pdfauthor={Hector Alexander Roman Potosi}
}

\pagestyle{fancy}
\fancyhf{}
\fancyhead[L]{\small Actividad 12 -- Fundamentos de IA}
\fancyhead[R]{\small Maestría en Ciencia de Datos}
\fancyfoot[C]{\thepage}

\setlength{\headheight}{15pt}
\setlength{\parskip}{5pt}
\setlength{\parindent}{0pt}

\setlist[itemize]{leftmargin=*,itemsep=2pt}

\lstset{
    basicstyle=\ttfamily\footnotesize,
    breaklines=true,
    frame=single,
    backgroundcolor=\color{gray!6},
    showstringspaces=false,
    columns=fullflexible
}

\newcommand{\repositorio}{
https://github.com/alexroma1126-design/actividad-12-chatbot-ia
}

\begin{document}

% ======================================================
% PORTADA
% ======================================================

\begin{titlepage}
\thispagestyle{empty}
\begin{center}

\vspace*{1cm}

{\color{azul}\rule{\textwidth}{1.2pt}}\\[0.8cm]

{\Large\bfseries
Universidad de Investigación de Tecnología Experimental Yachay
}\\[0.35cm]

{\large Maestría en Ciencia de Datos}\\[1.5cm]

{\color{azul}\LARGE\bfseries INFORME TÉCNICO}\\[0.45cm]

{\Huge\bfseries Actividad 12}\\[0.45cm]

{\Large
Chatbot con GPT-OSS, GroqCloud, Streamlit
y análisis de archivos CSV
}\\[0.8cm]

{\color{azul}\rule{0.78\textwidth}{0.8pt}}\\[1.5cm]

\begin{tabular}{@{}rl@{}}
\textbf{Asignatura:} & Fundamentos de Inteligencia Artificial\\
\textbf{Autor:} & Héctor Alexander Román Potosí\\
\textbf{Docente:} & Manuel Eugenio Morocho Cayamcela, PhD\\
\textbf{Fecha:} & 15 de noviembre de 2026\\
\textbf{Rama revisada:} & \texttt{main}\\
\textbf{Commit revisado:} & \texttt{4522728}\\
\end{tabular}

\vfill

\textbf{Repositorio público}\\[3mm]

\href{
https://github.com/alexroma1126-design/actividad-12-chatbot-ia
}{
\texttt{github.com/alexroma1126-design/actividad-12-chatbot-ia}
}\\[0.8cm]

{\small\color{gris}
Informe basado en el código publicado y
en la consigna del notebook de clase.
}

\end{center}
\end{titlepage}

\tableofcontents
\newpage

% ======================================================
\section{Resumen}
% ======================================================

Se desarrolló una aplicación web educativa con dos modos
de funcionamiento: conversación general en español y
consultas en lenguaje natural sobre archivos CSV.

La interfaz se implementó utilizando \textbf{Streamlit},
mientras que el procesamiento de datos se realizó mediante
\textbf{Pandas} y \textbf{SQLite en memoria}.

Para interpretar preguntas se empleó el modelo abierto
\texttt{openai/gpt-oss-20b}, desarrollado por OpenAI y
servido mediante \textbf{GroqCloud}.

La comunicación con el modelo utiliza el SDK de OpenAI
configurado con una URL compatible de Groq. Por tanto,
no es necesario consumir créditos de la API comercial
de OpenAI para realizar esta práctica.

El servicio gratuito de GroqCloud está sujeto a límites
de solicitudes y procesamiento.

El proyecto dispone de un repositorio público en GitHub
que contiene el código fuente, documentación,
un dataset de ejemplo y pruebas automatizadas.

El despliegue público en Streamlit Community Cloud
continúa pendiente de verificación.

% ======================================================
\section{Objetivos y relación con la consigna}
% ======================================================

La actividad propuesta en el notebook de la asignatura
se denomina:

\emph{Building a Simple Chatbot with Streamlit and
Hugging Face or OpenAI Models}.

La consigna establece la creación de una aplicación
que acepte preguntas, permita cargar archivos CSV,
responda sobre los datos y rechace consultas
ajenas al conjunto cargado.

Asimismo, solicita entregar el archivo Python
y el enlace público de la aplicación desplegada
en Streamlit Cloud, además de presentar los resultados
brevemente durante la clase.

\subsection{Objetivo general}

Desarrollar un chatbot interactivo que integre
inteligencia artificial generativa y cálculos
reproducibles para responder preguntas sobre
conjuntos de datos CSV desde una interfaz web.

\subsection{Objetivos específicos}

\begin{itemize}
\item Desarrollar una interfaz conversacional con Streamlit.
\item Integrar GPT-OSS mediante GroqCloud.
\item Transformar preguntas en consultas de solo lectura.
\item Analizar registros, columnas y estadísticas básicas.
\item Proteger las credenciales privadas del proveedor.
\item Documentar pruebas, resultados y limitaciones.
\item Publicar el código mediante GitHub.
\end{itemize}

% ======================================================
\section{Arquitectura de la solución}
% ======================================================

\subsection{Componentes tecnológicos}

\begin{table}[H]
\centering
\small
\caption{Componentes verificados en el repositorio.}

\begin{tabularx}{\textwidth}{@{}p{3cm}X@{}}
\toprule
\textbf{Componente} & \textbf{Responsabilidad}\\
\midrule

Python y Streamlit &
Interfaz web, conversación, carga y vista previa de archivos.\\

GPT-OSS 20B &
Comprensión del lenguaje natural y planificación de SQL.\\

GroqCloud &
Servicio de inferencia remota con cuota gratuita limitada.\\

SDK \texttt{openai} &
Cliente compatible con la API de Groq.\\

Pandas &
Lectura del CSV y operaciones descriptivas.\\

SQLite &
Consultas sobre tablas temporales en memoria.\\

\texttt{unittest} &
Pruebas de carga, cálculos, restricciones y errores.\\

GitHub &
Versionado y publicación del código fuente.\\

\bottomrule
\end{tabularx}
\end{table}

\subsection{Flujo de procesamiento}

\begin{center}
\fbox{
\begin{minipage}{0.91\textwidth}
\centering

\textbf{Usuario: pregunta o archivo CSV}\\[2mm]

$\downarrow$\\

\textbf{Interfaz Streamlit}\\[2mm]

\textbf{Chat general:}\\
GPT-OSS $\rightarrow$ respuesta en español\\[3mm]

\textbf{Modo CSV:}\\
Pandas carga el archivo y describe su esquema\\[2mm]

$\downarrow$\\

Preguntas simples $\rightarrow$ Pandas\\

Preguntas analíticas $\rightarrow$ GPT-OSS propone SQL\\[2mm]

$\downarrow$\\

\textbf{SQLite calcula el resultado}\\

$\downarrow$\\

\textbf{Streamlit presenta la respuesta y la tabla}

\end{minipage}
}
\end{center}

La separación entre interpretación lingüística
y cálculo numérico permite obtener resultados
que pueden contrastarse directamente con los datos.

Sin embargo, el modelo puede generar una consulta
SQL incorrecta. Por ello, los resultados deben
validarse frente al dataset original.

\subsection{Estructura del repositorio}

La organización del proyecto publicado es:

\begin{lstlisting}
actividad-12-chatbot-ia/
  app.py
  README.md
  requirements.txt
  .gitignore
  .streamlit/secrets.toml.example
  datos/clientes_ejemplo.csv
  tests/test_app.py
\end{lstlisting}

% ======================================================
\section{Implementación y seguridad}
% ======================================================

La aplicación dispone de dos modalidades principales.

\subsection{Chat general}

Permite ingresar preguntas abiertas en español
y recibir respuestas generadas por GPT-OSS.

El historial de conversación se mantiene durante
la sesión de Streamlit.

\subsection{Preguntas sobre un CSV}

El usuario carga un archivo CSV y la aplicación
muestra una vista previa de sus registros.

El archivo se procesa mediante Pandas.

La aplicación puede responder directamente preguntas
sobre cantidad de registros, columnas y valores
mínimos o máximos.

Para consultas analíticas, el modelo interpreta
la pregunta y propone una sentencia SQL que
posteriormente se ejecuta mediante SQLite.

El límite de tamaño establecido para archivos
es de \textbf{3 MB}.

La cantidad máxima de filas mostradas por consulta
es de \textbf{100 registros}.

\subsection{Restricciones de seguridad}

El código exige consultas SQL de tipo
\texttt{SELECT} y utiliza un autorizador de SQLite
para limitar las operaciones permitidas.

Las consultas se ejecutan sobre la tabla temporal
\texttt{datos}, evitando modificaciones a
archivos externos.

La credencial privada se recupera desde la variable
\texttt{GROQ\_API\_KEY} o los secretos de Streamlit.

El archivo privado \texttt{.streamlit/secrets.toml}
se encuentra excluido mediante \texttt{.gitignore}
y no aparece en el árbol publicado del repositorio.

Estas medidas son apropiadas para una demostración
educativa, aunque no reemplazan una auditoría
integral de seguridad.

El esquema y algunas categorías del CSV se envían
al proveedor para interpretar consultas.

Por ello, no deben utilizarse datos personales
ni información sensible sin la autorización
correspondiente.

% ======================================================
\section{Datos de prueba y verificación}
% ======================================================

\subsection{Dataset de ejemplo publicado}

El archivo \texttt{datos/clientes\_ejemplo.csv}
incluido en GitHub contiene \textbf{20 registros}
y cinco columnas:

\begin{itemize}
\item \texttt{ClienteID}
\item \texttt{Genero}
\item \texttt{Edad}
\item \texttt{IngresoAnualUSD}
\item \texttt{PuntuacionGasto}
\end{itemize}

Los resultados de referencia se presentan
en la Tabla siguiente.

\begin{table}[H]
\centering
\small
\caption{Resultados verificables del CSV de ejemplo.}

\begin{tabularx}{\textwidth}{@{}Xr@{}}
\toprule
\textbf{Indicador} & \textbf{Resultado}\\
\midrule
Número de registros & 20\\
Número de columnas & 5\\
Registros con género F & 10\\
Registros con género M & 10\\
Promedio de edad & 33,45\\
Edad mínima & 19\\
Edad máxima & 51\\
Ingreso anual máximo (USD) & 48\,000\\
\bottomrule
\end{tabularx}
\end{table}

Durante las pruebas locales también se utilizó
otro dataset relacionado con diabetes,
con \textbf{768 registros y 9 columnas}.

Ese dataset corresponde a una prueba distinta
y no es el archivo de clientes publicado
en el repositorio de GitHub.

\subsection{Pruebas automatizadas}

El archivo \texttt{tests/test\_app.py}
contiene \textbf{11 métodos de prueba} relacionados con:

\begin{itemize}
\item Lectura y validación de archivos CSV.
\item Conteo de filas.
\item Agrupaciones.
\item Cálculo de promedios.
\item Rechazo de instrucciones SQL no autorizadas.
\item Interpretación simulada de preguntas fuera de contexto.
\item Manejo de errores de autenticación y cuotas.
\end{itemize}

Algunas pruebas utilizan un cliente ficticio
y no realizan solicitudes reales a Groq.

La revisión del repositorio confirma la existencia
de estas pruebas, pero no equivale a una ejecución
reciente de integración continua.

\subsection{Matriz de aceptación}

\begin{table}[H]
\centering
\small
\caption{Casos recomendados para la demostración.}

\begin{tabularx}{\textwidth}{@{}p{5.3cm}X@{}}
\toprule
\textbf{Pregunta o prueba} &
\textbf{Resultado esperado}\\
\midrule

Saludo en chat general &
Respuesta en español mediante GPT-OSS.\\

¿Cuántos registros tiene el CSV? &
20 registros.\\

¿Cuáles son las columnas? &
Listado de cinco encabezados.\\

¿Cuál es el promedio de Edad? &
33,45 o su redondeo explícito.\\

¿Cuántos registros hay por género? &
F: 10; M: 10.\\

¿Cuál es la capital de Francia? (modo CSV) &
Rechazo por estar fuera del dataset.\\

Intento de SQL de escritura &
Operación bloqueada.\\

\bottomrule
\end{tabularx}
\end{table}

% ======================================================
\section{Limitaciones y trabajo pendiente}
% ======================================================

\begin{itemize}

\item El modelo puede rechazar consultas válidas
o producir SQL incorrecto.

\item Las definiciones de columnas son orientativas
y no sustituyen el diccionario original de datos.

\item El plan gratuito de Groq está sujeto
a límites de uso.

\item Una aplicación pública puede recibir solicitudes
que consuman la cuota disponible.

\item Las pruebas publicadas todavía no cubren
todas las variantes de preguntas sobre columnas
y extremos numéricos.

\item Falta verificar el despliegue público en
Streamlit Community Cloud.

\item Es recomendable adjuntar capturas de las
pruebas finales sin exponer credenciales.

\end{itemize}

% ======================================================
\section{Reproducibilidad y enlaces}
% ======================================================

\subsection{Ejecución local}

Desde la carpeta raíz del proyecto se utilizan
los siguientes comandos:

\begin{lstlisting}
python -m pip install -r requirements.txt
python -m unittest discover -s tests -v
python -m streamlit run app.py
\end{lstlisting}

La aplicación se abre localmente en:

\url{http://localhost:8501}

La clave privada se configura mediante
\texttt{GROQ\_API\_KEY}.

Nunca debe incorporarse al informe,
al repositorio público ni a las capturas.

\subsection{Despliegue en Streamlit Cloud}

Para completar la publicación se debe:

\begin{enumerate}
\item Conectar el repositorio de GitHub.
\item Seleccionar la rama \texttt{main}.
\item Establecer \texttt{app.py} como archivo principal.
\item Configurar la clave privada en los secretos de Streamlit.
\item Publicar la aplicación.
\item Comprobar ambos modos de funcionamiento.
\end{enumerate}

La plataforma de despliegue es:

\url{https://share.streamlit.io}

\subsection{Enlaces del proyecto}

\begin{tabularx}{\textwidth}{@{}p{3cm}X@{}}
\toprule
\textbf{Recurso} & \textbf{Enlace o estado}\\
\midrule

Repositorio GitHub &
\url{https://github.com/alexroma1126-design/actividad-12-chatbot-ia}\\

Código Python &
\url{https://github.com/alexroma1126-design/actividad-12-chatbot-ia/blob/main/app.py}\\

Commit revisado &
\texttt{4522728} (rama \texttt{main})\\

App pública Streamlit &
\textbf{Pendiente de desplegar o verificar.}\\

\bottomrule
\end{tabularx}

% ======================================================
\section{Conclusiones}
% ======================================================

Se implementó un chatbot híbrido que combina
inteligencia artificial generativa y procesamiento
tabular mediante herramientas de Python.

GPT-OSS interpreta las preguntas, mientras
que Pandas y SQLite permiten realizar operaciones
sobre los datos cargados.

Esta arquitectura mejora la verificabilidad
de los resultados frente a un chatbot
que responde utilizando únicamente texto generado.

El proyecto dispone de código fuente público,
documentación, pruebas definidas y mecanismos
básicos de protección de credenciales.

GroqCloud permitió ejecutar un modelo abierto
de OpenAI sin consumir créditos de la API comercial
de OpenAI, aunque mantiene restricciones de uso.

Para cumplir completamente los entregables
de la asignatura, es indispensable verificar
la publicación en Streamlit Community Cloud
y entregar su enlace junto con \texttt{app.py}
mediante Moodle.

% ======================================================
\section*{Referencias}
\addcontentsline{toc}{section}{Referencias}
% ======================================================

\begin{enumerate}

\item Morocho Cayamcela, M. E. (2026).
\emph{OpenAI API}.
Notebook de Fundamentos de Inteligencia Artificial,
actividad final.

\item Román Potosí, H. A. (2026).
\emph{Actividad 12: Chatbot GPT-OSS con Groq y Streamlit}.
\url{https://github.com/alexroma1126-design/actividad-12-chatbot-ia}.

\item Groq.
\emph{Documentación del modelo GPT-OSS 20B}.
\url{https://console.groq.com/docs/model/openai/gpt-oss-20b}.

\item Streamlit.
\emph{Community Cloud: desplegar aplicaciones}.
\url{https://docs.streamlit.io/deploy/streamlit-community-cloud}.

\end{enumerate}

\end{document}